\section{Related Work}
\label{sec:related}

\paragraph{Shapley-value feature attribution and its estimators.}
SHAP casts feature attribution as an additive Shapley-value problem \citep{lundberg2017unified}; exact computation is exponential in the number of features, so model-agnostic explainers estimate it by sampling coalitions (Kernel SHAP) or feature permutations. The variance of these estimators has been analysed and reduced by paired sampling and regression-based variants \citep{covert2021improving}, by better permutation sampling \citep{mitchell2022sampling}, and by leverage-score sampling with error guarantees \citep{musco2025leverage}. Sampling estimators of Shapley-type contributions go back to \citet{strumbelj2014explaining}, and tree ensembles admit deterministic algorithms \citep{lundberg2020local}. These works study estimator error as a function of the number of samples; we ask the deployment question that follows from them: given the budget a practitioner actually uses, \emph{is one run enough to treat the attribution as the explanation of that prediction}, and for which instances and models.

\paragraph{What the target is.}
The quantity SHAP estimates depends on how the value function takes expectations over absent features and on the background data \citep{chen2020true}; background-set size already affects explanation stability for neural models \citep{yuan2023background}. We therefore define error and agreement relative to a fixed model, value function and background set, and treat background resampling as a separate source of variation.

\paragraph{Disagreement and multiplicity.}
Explanations can disagree across methods \citep{krishna2024disagreement}, and equally accurate models can disagree on individual predictions \citep{marx2020predictive}. We use the second as a comparison standard (explainer-run noise is judged against between-refit variation instead of in isolation); between-method disagreement is not evaluated here.

\paragraph{Instability of sampling-based explainers.}
Other randomised explainers are known to vary between runs: for LIME, repeated runs can select different features, which motivated stability indices computed from repeated explanations \citep{visani2022stability}, and robustness metrics for interpretability methods in general were proposed by \citet{alvarezmelis2018robustness}. These works take repeatability as the criterion; for SHAP estimators we add comparison with an exact target and with other sources of variation.

\paragraph{Stability and agreement measures.}
Set-stability of selected features \citep{nogueira2018stability} and rank-biased overlap for top-weighted, incomplete rankings \citep{webber2010rbo} give chance-aware ways to compare feature sets and rankings across runs; we adopt them alongside error and sign-agreement measures.

\paragraph{Positioning.}
The estimator-variance literature above gives error as a function of the number of samples for a given algorithm; the stability and multiplicity literature compares explanations across models or methods. The error-versus-samples curves of Figure~\ref{fig:budget} are in the spirit of \citet{covert2021improving}, \citet{mitchell2022sampling} and \citet{musco2025leverage}; what we add is their reading at the library's actual budgets for the question of top-feature recovery, together with the other notions below. The metrics themselves are not new; the combination is applied to one deployment question. We separate repeatability, accuracy against an exact target, resolvability of the ordering and noise relative to background and refit variation, evaluate them at default and neighbouring budgets of the most used implementation, and state the conditions under which the single-run practice holds. We read only the open literature cited here; other related work may exist that we did not consult.
